\chapter{Benchmark and Reproducibility Protocol}
\label{ch:benchmark-reproducibility-protocol}

\section{Why ORIUS-Bench is central to the thesis}
Once the runtime object has been fixed, the next question is whether it can be measured in a
way that matches the theory rather than only the software stack. ORIUS-Bench is the answer
to that question. It is not infrastructure attached after the science; it is the
measurement contract that makes the dissertation's central object observable. If OASG is the
gap between observed-state legality and true-state safety, then the benchmark must expose
both surfaces explicitly. A benchmark that measures only forecast error, controller cost, or
observed-state success cannot validate the ORIUS claim because it cannot see the hidden
failure event the book is about.

\section{Dual-trajectory evaluation}
Every ORIUS-Bench episode therefore carries two trajectories:
\begin{itemize}
  \item the \emph{observed} trajectory available to the runtime after degradation; and
  \item the \emph{true} trajectory used only by the replay or evaluation harness to judge
  whether release was physically safe.
\end{itemize}
That separation makes the benchmark stricter than a conventional controller comparison. A
nominal controller may still appear coherent on the observed state while the replay harness
records a true-state violation. ORIUS-Bench treats that divergence as first-class evidence.

\section{Deterministic replay and fault schedules}
The benchmark uses deterministic fault schedules so that comparisons are reproducible and so
that certificate traces can be replayed under the same degraded-observation regime. Fault
families include dropout, stale holds, delay, reordering, spikes, drift, and domain-local
analogs. The scientific point is not merely that faults exist. It is that the same fault
grammar can be replayed against one release contract across domains.

Strict replay discipline matters for another reason: it prevents the dissertation from
hiding behind one-off operational anecdotes. A run must be reproducible through a manifest,
fault schedule, and artifact family. If a result cannot be regenerated from the tracked
surface, it is not part of the defended thesis claim.

\section{Metrics that match the runtime object}
ORIUS-Bench measures the runtime layer through seven core metrics:
\begin{description}
  \item[TSVR] true-state violation rate, the primary safety outcome;
  \item[OASG] the observed-safe / true-unsafe discrepancy rate;
  \item[CVA] certificate-valid action rate or validity surface;
  \item[GDQ] graceful degradation quality under fallback or reduced-operation modes;
  \item[IR] intervention rate, capturing how often the runtime must repair or deny release;
  \item[AC] audit completeness, measuring whether release is reconstructable after the fact;
  \item[RL] runtime latency, measuring whether the protection layer remains operationally
  usable.
\end{description}

This metric family matters because it keeps the book from collapsing into a single
leaderboard score. ORIUS is not about RMSE alone, and it is not about controller cost alone.
It is about safe release under degraded observation. The benchmark therefore measures
whether release remained safe, whether intervention was required, whether the certificate
remained meaningful, and whether the runtime stayed fast enough to matter.

\IfFileExists{paper/assets/tables/generated/cross_domain_oasg_table.tex}{%
\begin{table}[htbp]
\centering
\caption{Cross-Domain OASG (Observed-Action Safety Gap) Rates.}
\label{tbl:cross_domain_oasg}
\begin{tabular}{lccc}
\toprule
\textbf{Domain} & \textbf{Baseline OASG} & \textbf{ORIUS OASG} & \textbf{Reduction} \\
\midrule
Battery (Ref.) & 0.0083 & 0.0000 & 100.0\% \\
Healthcare & 0.1945 & 0.0000 & 100.0\% \\
Vehicle (AV) & 0.2893 & 0.0002 & 99.9\% \\
\bottomrule
\end{tabular}
\end{table}%
}{%
\begin{table}[htbp]
\centering
\caption{Cross-Domain OASG (Observed-Action Safety Gap) Rates.}
\label{tbl:cross_domain_oasg}
\begin{tabular}{lccc}
\toprule
\textbf{Domain} & \textbf{Baseline OASG} & \textbf{ORIUS OASG} & \textbf{Reduction} \\
\midrule
Battery (Ref.) & 0.0083 & 0.0000 & 100.0\% \\
Healthcare & 0.1945 & 0.0000 & 100.0\% \\
Vehicle (AV) & 0.2893 & 0.0002 & 99.9\% \\
\bottomrule
\end{tabular}
\end{table}%
}

\section{Strict universal mode and legacy compatibility}
The dissertation adopts a strict universal reading of the benchmark. Domain-neutral
violation predicates, observed-state legality predicates, and constraint margins must flow
through the adapter contract. Legacy battery-shaped compatibility fields may remain in the
codebase for historical or bounded-support reasons, but they are not allowed to define the
universal benchmark story in the manuscript.

That distinction keeps the witness row from quietly redefining the universal object. Battery
is the deepest defended row because its evidence chain is strongest, not because the
benchmark is secretly battery-specific. ORIUS-Bench must therefore be written in
domain-agnostic terms even when the strongest current replay surface is battery-backed.

\section{Artifact provenance and claim locking}
The benchmark protocol is inseparable from the claim-lock pipeline. The canonical submission
surface is \texttt{orius\_book.tex}, while the mirrored internal controller remains
\texttt{paper/paper.tex}; their headline tables and figures are backed by
manifests, publication matrices, and governed artifact paths. This is what prevents the
manuscript from becoming a persuasive story about results that cannot later be checked.

The repository's publication surfaces already enforce this discipline. The parity matrix,
submission scorecard, runtime/governance matrices, and deployment-validation scope are not
mere reporting layers. They are the machine-readable surfaces that keep prose aligned with
evidence status.

\section{Reproducibility checklist}
For a result to count as defended in this dissertation, the following conditions must hold:
\begin{enumerate}
  \item the domain row must have an explicit safety predicate and adapter semantics;
  \item the replay must expose true and observed trajectories;
  \item the fault schedule and manifests must be reproducible;
  \item the result must land in a tracked artifact family; and
  \item the claim validator must allow the corresponding manuscript language.
\end{enumerate}
This checklist is stricter than ``the code ran once.'' It is the mechanism by which ORIUS
turns software artifacts into scientific evidence.

\section{Benchmark takeaway}
What this chapter adds to the argument is a measurement discipline equal to the runtime
object introduced in Chapters~3 and~4. ORIUS-Bench makes degraded-observation safety
measurable through dual-trajectory replay, deterministic fault injection, runtime metrics
matched to the release object, and claim-locked provenance. Equal-domain maturity still
depends on which rows clear those requirements, but the measurement contract itself is now
stable across the dissertation.
